%قبل از شروع مسئله ترجیحاً چیزی ننویسید. اما اگر خواستید بنویسید، در قالب لیستی (فقرات و شمارش) بنویسید.
%برای اضافه کردن نمره، بعدی از شروع‌مسئله، [x] را اضافه کنید و به جای x  نمره مورد نظر را قرار دهید. برای زیرمسئله‌ها هم به همین صورت است.
%\شروع‌مسئله[10]
\شروع‌مسئله [20]
\متن‌سیاه{بلاک‌چین.}
در این سوال شما باید سیستم ساده شده‌ای از بلاک‌چین یا زنجیرهٔ بلوکی را پیاده‌سازی کنید. ابتدا برای آشنایی با بلاک‌چین پیشنهاد می‌شود از 
\href{https://blog.goodaudience.com/blockchain-for-beginners-what-is-blockchain-519db8c6677a}{این لینک}
 استفاده کنید. 
 سرآیند هر بلوک ۸۰ کاراکتر اول آن را تشکیل می‌دهد و شامل شش فیلد زیر می‌باشد.
\begin{center}
	\begin{tabular}{|l|r|c|}
		\hline
		\multicolumn{1}{|c|}{نام فیلد}
		            & \multicolumn{1}{c|}{کاربرد}               
		            	          & 
		            	          طول (بایت) 
		            	          \\		\hline		\hline
		Version        & نسخهٔ بلوک                        & ۴          \\		\hline
		hashPrevBlock  & درهم‌سازی ۲۵۶ بیتی بلوک قبلی     & ۳۲         \\		\hline
		hashMerkleRoot & درهم‌سازی ۲۵۶ بیتی داده‌های بلوک & ۳۲         \\		\hline
		Time           & زمان حال به ثانیه                & ۴          \\		\hline
		Bits           & میزان سختی تولید بلاک            & ۴          \\		\hline
		Nonce          & یک عدد دلخواه                    & ۴         
\\		\hline
	\end{tabular}
\end{center}
بلوک‌ها یک لیست پیوندی را تشکیل می‌دهند که بلاک‌چین یا زنجیرهٔ بلوکی نام دارد. هر بلوک این زنجیره با فیلد hashPrevBlock به بلوک قبلی خود متصل می‌شود. در این سؤال روی یک بلوک از زنجیره تمرکز می‌کنیم. بعد از جمع‌آوری اطلاعات اصلی یک بلوک، سرآیند آن تشکیل می‌شود که ۵ فیلد اول آن مشخص است. برای این که بلاک مجوز اضافه شدن به زنجیره را داشته باشد، باید فیلد ششم یا nonce به صورتی انتخاب شود که اگر تابع \متن‌لاتین{SHA-256} را \زیرخط{دوبار} روی سرآیند بلوک اجرا کنیم، به تعداد خاصی صفر اول آن ظاهر شود. این شرط به این منظور گذاشته شده که اضافه کردن بلوک جدید به زنجیره کار سختی باشد و با وجود تلاش جهانی، حدوداً هر ۱۰ دقیقه یک بلوک به زنجیره اضافه شود.
در ورودی به شما ۷۶ کاراکتر اول سرآیند داده می‌شود و وظیفۀ شما پیدا کردن nonce مناسب و تکمیل بلوک است.  به این کار ماینینگ گفته می‌شود. سرآیند در ورودی به صورت هگز (مبنای ۱۶) داده شده است و بنابراین به صورت ۱۵۲ کاراکتر ظاهر می‌شود. در خروجی نیز nonce را باید به صورت هگز بنویسید که ۴ کاراکتر آن به صورت ۸ کاراکتر هگز ظاهر می‌شود. بعد از اضافه کردن \متن‌لاتین{nonce}، درهم‌سازی سرآیند بلوک باید ۴ تا صفر در نمایش هگر خود داشته باشد (این عدد در واقعیت ۱۸ صفر می‌باشد و به تدریج ممکن است زیاد شود).

توجه کنید ورودی و خروجی به صورت هگز هستند ولی محاسبات باید روی اطلاعات اصلی (باینری) انجام شود. همچنین در تبدیل درهم‌سازی به هگز، ممکن است رشته به صورت برعکس تبدیل شود و صفرها در برخی توابع ابتدا و یا انتهای رشته قرار گیرند. برای همین ابتدا برنامه خود را روی ورودی نمونه تست کنید. شما می‌توانید از کتابخانهٔ hashlib پایتون استفاده کنید.
در
\href{https://webbtc.com/block/000000000000000001f942eb4bfa0aeccb6a14c268f4c72d5fff17270da771b9}{این لینک}
بلوک کامل ورودی را می‌توانید پیدا کنید (\متن‌لاتین{nonce} اصلی این بلوک با خروجی نمونه متفاوت است).

\begin{center}\small
	\begin{tabular}{|c|c|}
		\hline
		خروجی                 
		            & ورودی    \\ \hline
\texttt{7735c84e} &
\parbox[t]{8.1cm}{\normalsize\raggedright{\texttt{0200000066191da95594aeda1a98a19ff054a88a\\
510754e2a4d93e0a00000000000000008485ae79\\
7312b2cb37dfb1aac11d7c5ad9dd84364bbe26ff\\
a781853996587d9b10a06555f5861618}}\vspace{0.2cm}}
%\lr{}
 \\ \hline
	\end{tabular}
\end{center}
درهم‌سازی ورودی+خروجی نمونه:
{\small\texttt{0000a22ba24c4bae0cfc072d79b29213364a1b339c1a78fd927d0d3a511231a0}}



%اگر مسئله چند قسمت داشت به این صورت عمل کنید:
%\شروع{زیرمسائل}
%\زیرمسئله[0] % مسئلۀ ۱آ
%\زیرمسئله % مسئلۀ ۱ب
%\زیرمسئله % مسئلۀ ۱پ
%\زیرمسئله % مسئلۀ ۱ت
%\پایان{زیرمسائل}

\solution{
در این سؤال کافی است در یک حلقه \متن‌لاتین{nonce}های مختلف را بررسی کنید تا جواب پیدا شود. در عین سادگی راه‌حل، این سؤال، الگوریتم استفاده شده در \متن‌لاتین{miner}های بلاک‌چین را نشان می‌دهد. در واقع اجرای همین الگوریتم در سال بیش از ۶۶ تراوات‌ساعت برق مصرف می‌کند که بیشتر از مصرف برق بسیاری از کشورها مثل پرتغال است. لازم به ذکر است \متن‌لاتین{miner}ها برای اجرای بهتر و سریع‌تر این الگوریتم که منجر به کشف بیت‌کوین و کسب ثروت می‌شود، سخت‌افزارهای مخصوصی طراحی و تولید کرده‌اند.

فایل \متن‌لاتین{\متن‌تایپ{blockchain\_sol.py}} شامل کد راه‌حل به عنوان نمونه ضمیمه شده است. 

}




%%%%%%%%%%%%%%%%%%%%%%%%%%%%%%%%%%%%%%%%%%%%%%%%%%%%%
% برخی دستورات مفید برای تمرین‌ها. %
%%%%%%%%%%%%%%%%%%%%%%%%%%%%%%%%%%%%%%%%%%%%%%%%%%%%%

% فرمول‌های نمونه:
%$f(x) = \Theta(x)$
%$T(x, y) \leq \log(x) + 2^y + \binom{2n}{n}$
% {\tt code\_function}


% لیست موارد بدون شماره‌گذاری:
%\شروع{فقرات}
%    \فقره مورد اول 
%        \شروع{فقرات}
%            \فقره مورد اول داخلی 
%                \شروع{فقرات}
%                    \فقره اولین مورد خیلی داخلی 
%                    \فقره دومین مورد خیلی داخلی 
%                \پایان{فقرات}
%            \فقره مورد دوم داخلی 
%        \پایان{فقرات}
%    \فقره مورد دوم 
%\پایان{فقرات}

% لیست موارد شماره‌گذاری شده:
%\شروع{شمارش}
%    \فقره یک مورد 
%    \فقره یک مورد دیگر 
%    \فقره باز هم یک مورد دیگر
%\پایان{شمارش}

% برای نوشتن یک فرمول چند خطی که یک شماره بخوره. برای حالت‌های دیگر به سایت زیر مراجعه کنید:
% https://www.overleaf.com/learn/latex/Aligning_equations_with_amsmath
%\begin{align} 
%    \begin{split}
%        (x+y)^3 &= (x+y)^2(x+y) \\
%                &= (x^2+2xy+y^2)(x+y) \\
%                &= (x^3+2x^2y+xy^2) + (x^2y+2xy^2+y^3) \\
%                &= x^3+3x^2y+3xy^2+y^3
%    \end{split}                                 
%\end{align}

% نوشتن ماتریس و آرایه دوبعدی
%\begin{align}
%    A = 
%    \begin{bmatrix}
%        A_{11} & A_{21} \\
%        A_{21} & A_{22}
%    \end{bmatrix}
%\end{align}

% اضافه کردن تصویر با شرح. برای اضافه کردن تصویر بدون شرح، فقط دستور درج‌تصویر را استفاده کنید.
%\شروع{شکل}[h!]
%\تنظیم‌ازوسط
%  \درج‌تصویر[پهنا=0.5\پهنای‌متن]{img.jpg}
%\شرح{شرح تصویر}
%\برچسب{تصویر:۱}
%\پایان{شکل}
%\پایان{زیرمسائل}

%کد زیر را ببینید.
%%\lstinputlisting{source_filename.py}
%\begin{latin}
%\begin{lstlisting}
%# Import the modules
%import sys
%import random
%
%ans = True
%
%while ans:
%question = raw_input("Ask the magic 8 ball a question: (press enter to quit) ")
%
%answers = random.randint(1,8)
%
%if question == "":
%sys.exit()
%
%elif answers == 1:
%print "It is certain"
%
%elif answers == 2:
%print "Outlook good"
%
%elif answers == 3:
%print "You may rely on it"
%
%elif answers == 4:
%print "Ask again later"
%
%elif answers == 5:
%print "Concentrate and ask again"
%
%elif answers == 6:
%print "Reply hazy, try again"
%
%elif answers == 7:
%print "My reply is no"
%
%elif answers == 8:
%print "My sources say no"
%\end{lstlisting}
%\end{latin}
